\section{April 09}


  Let $(M, \omega)$ be a symplectic manifold and $G$ a Lie group acting on $M$ by symplectomorphisms.
  Let $\frakg$ be the Lie algebra of $G$ and $\frakg^*$ its dual.
  Assume the action is Hamiltonian and $\mu: M \to \frakg^*$ is the moment map.

  % \begin{remark}
  Let $f,h \in C^\infty(\frakg^*)$.
  If $\frakg$ is abelian, then 
  \[ \{f,h\}(\xi) = \langle \xi, [df_\xi, dh_\xi] \rangle = 0. \]
  In this case, the Poisson structure on $\frakg^*$ is trivial.
  Then we have $\{f \circ \mu, h \circ \mu\} = 0$.
  % \end{remark}

  If $\dim \frakg = \dim M/2$ and the action of $G$ on $M$ having discrete kernel, then we can obtain an integrable system on $M$ by pulling back functions on $\frakg^*$ via the moment map $\mu$.


  \begin{lemma}
    For $p \in M$, $Y \in T_p M$ and $X \in \frakg$, we have 
    \[ \bigl(d \mu_p(Y)\bigr)(X) = \omega_p(\tilde{X}, Y). \]
  \end{lemma}
  \begin{proof}
    Let $\gamma$ be a curve in $M$ such that $\gamma(0) = p$ and $\gamma'(0) = Y$.
    Then we have
    \begin{align*}
      \bigl(d \mu_p(Y)\bigr)(X) &= \frac{d}{dt}\Big|_{t=0} \mu(\gamma(t))(X) = \frac{d}{dt}\Big|_{t=0} \langle \mu(\gamma(t)), X \rangle \\
        & = \frac{d}{dt}\Big|_{t=0} \tilde{\mu}(X)(\gamma(t)) = Y(\tilde{\mu}(X)) \\
        & = d \tilde{\mu}(X) (Y) = -i_{\tilde{X}} \omega (Y) = \omega_p(\tilde{X}, Y).
    \end{align*} 
    We are done.
  \end{proof}

  Recall we have the co-adjoint action of $G$ on $\frakg^*$.

  \begin{lemma}
    Assume $G$ is connected.
    The moment map $\mu$ is $G$-equivariant with respect to the given action of $G$ on $M$ and the co-adjoint action of $G$ on $\frakg^*$.
  \end{lemma}
  \begin{proof} 
    Since $G$ is generated by a neighborhood of the identity, we only need to check the condition for the image of the exponential map $\exp: \frakg \to G$.
    WTS
    \begin{equation}
      \mu\left( \exp(tX) \cdot p \right) = \Ad^*_{\exp(tX)} \mu(p), \quad \forall t \in \bbR.
    \end{equation} 
    Taking derivative, it is equivalent to show that
    \[ d \mu_p(\tilde{X}_M) = \tilde{X}_{\frakg^*}. \]

    For $Y \in \frakg$, we have
    \begin{align*}
      \bigl(d \mu_p(\tilde{X}_M)\bigr)(Y) & = \omega_p(\tilde{X}_M, \tilde{Y}_M) = \{ \tilde{mu}(X), \tilde{\mu}(Y) \} (p) \\
        & = - \tilde{\mu}([X,Y])(p) = \langle \mu(p), -[X,Y] \rangle \\
        & = \langle \mu(p), \frac{d}{dt}\Big|_{t=0} \Ad_{\exp(-tX)} Y \rangle  = \frac{d}{dt}\Big|_{t=0} \langle \Ad^*_{\exp(tX)} \mu(p), Y \rangle \\
        & = \bigl( \tilde{X}_{\frakg^*} \bigr)(Y).
    \end{align*}
    We are done.
  \end{proof}

  \begin{example}
    Let $G = \SO(3)$ and $\frakg = \so(3)$.
    We can identify $\so(3)$ with $\bbR^3$ via the isomorphism
    \[ \Phi: \begin{pmatrix}
        0 & -z & y \\
        z & 0 & -x \\
        -y & x & 0
      \end{pmatrix} \mapsto (x,y,z). \]
    View $\bbR^3$ as a Lie algebra with the cross product as the Lie bracket.
    Then $\Phi$ is a Lie algebra isomorphism, i.e., $\Phi([A,B]) = \Phi(A) \times \Phi(B)$ for $A,B \in \so(3)$.
    Using the standard inner product on $\bbR^3$, we can identify $T\bbR^3 \cong T^* \bbR^3 \cong \bbR^3 \times \bbR^3$.
    Let $M = T^* \bbR^3$ and $\omega$ be the standard symplectic form on $T^* \bbR^3$.
    We have an action of $\SO(3)$ on $M$ given by $g \cdot (x,y) = (gx, gy)$ for $g \in \SO(3)$ and $(x,y) \in M$.

    This action is Hamiltonian and the moment map $\mu: M \to \so(3)^* \cong \bbR^3$ is given by $\mu(q,v) = q \times v$.
    In physics, view $q$ as the position and $v$ as the velocity, then $\mu(q,v)$ is the angular momentum map.
  \end{example}

  \begin{lemma}
    If $(M,\omega)$ is compact, then any Hamiltonian vector field has a zero.
  \end{lemma}
  \begin{proof}
    Let $X$ be a Hamiltonian vector field on $M$ with Hamiltonian function $f$.
    Since $M$ is compact, $f$ attains its maximum at some point $p \in M$.
    Then we have $df_p = 0$, which implies that $X_p = 0$.
  \end{proof}

  \begin{example}
    Let $G = S^1$ acting on $M = S^1 \times S^1$ by $g \cdot (x,y) = (x + g, y)$ for $g \in S^1$ and $(x,y) \in M$.
    Consider $X = 1 \in \bbR \cong \frakg$ and the corresponding vector field $\tilde{X}$ on $M$.
    Note that $\tilde{X}$ is a non-vanishing vector field on $M$.
    Then $\tilde{X}$ cannot be a Hamiltonian vector field by the above lemma.
    In particular, the action of $S^1$ on $M$ is not Hamiltonian.
  \end{example}

  \begin{example}
    Let $f_1,\cdots,f_k \in C^\infty(M)$ satisfy $\{f_i, f_j\} = 0$ for all $i,j$.
    Let $X_i$ be the Hamiltonian vector field corresponding to $f_i$.
    Assume $X_i$ are complete.
    Let $\phi_i^t$ be the flow of $X_i$.
    Consider the action of $G = \bbR^k$ on $M$ given by $(t_1,\cdots,t_k) \cdot p = \phi_1^{t_1} \circ \cdots \circ \phi_k^{t_k}(p)$ for $p \in M$ and $(t_1,\cdots,t_k) \in G$.
    This is well-defined since $X_i$ commute with each other.

    This action is Hamiltonian and the moment map $\mu: M \to \frakg \cong \bbR^k$ is given by $\mu(p) = (f_1(p),\cdots,f_k(p))$.
  \end{example}
